\chapter*{The software}
\addcontentsline{toc}{chapter}{The software}
\markboth{The software}{The software}

Every number in this book comes out of code, and the code is public. This
chapter says where it is, what is in it, and --- equation by equation --- which
routine evaluates each result and which test or script checks it.

A book of this kind accumulates a great many closed forms, and a reader who
wants to use one of them needs to know what it was checked against. Table~\ref{tab:map} answers that for every
numbered equation the book actually computes with.

\section*{Where it is}

One repository, holding two packages that install separately and carry their
own test suites:

\begin{quote}
\small\url{https://github.com/av2atgh/chygraph}
\end{quote}

\noindent
\code{percolation/} carries percolation on chygraphs: the self-consistency map,
the threshold tensor, the critical amplitude, correlated layers, and the
constructions of Chapters~\ref{ch:percolation}--\ref{ch:epidemics}. It is
symbolic throughout, on \code{sympy}.

\code{statmech/} carries everything from Chapter~\ref{ch:statmech} onward: the
cavity recursion with a general interior, the branching matrix, the Ising and
hitting-set and core-percolation solvers, the Bethe free energy, and the
region-graph calculation of Chapter~\ref{ch:gbp}. It depends on the first ---
\code{statmech.\allowbreak stability} imports \code{Percolation\allowbreak Matrix}
--- which is why the two live together.

A third package, \code{chygraph}, is one namespace over both, so a routine can
be reached without knowing which of the two defines it. Table~\ref{tab:map}
names the defining package throughout, since that is what a reader looking for
the source wants.

A third source is the book's own \code{figs/} directory, which holds thirteen
scripts covering Chapter~\ref{ch:data} onward, most of them named for the
chapter they serve. Those scripts generate the figures, and for
Chapters~\ref{ch:potts}, \ref{ch:colouring} and~\ref{ch:satisfiability} they
also carry the calculations themselves --- the last two having no manuscript
behind them at all.

Two further directories matter, because several results in the book are read
from them rather than recomputed. \code{statmech/\allowbreak examples/}
runs the comparisons a reader is most likely to want --- the clustering
comparison of Sec.~\ref{sec:clustering}, the hard-field against soft-field
comparison of Sec.~\ref{sec:soft}. \code{statmech/\allowbreak probe/} holds the
expensive measurements, with their outputs cached under
\code{statmech/\allowbreak probe/\allowbreak results/}: the Monte
Carlo of Eq.~\eqref{eq:mc}, the leaf-removal validation of
Sec.~\ref{sec:cover-checks}, the hyperbolic-graph scan behind
Fig.~\ref{fig:hrg}, the sixteen real networks of Table~\ref{tab:realcore}, and
the seven measurements of Part~IV --- the cavity method on
the clique chygraph behind Fig.~\ref{fig:cavityfail}, the sixty region graphs
behind Fig.~\ref{fig:gbp}, the placed-motif instances of
Sec.~\ref{sec:karrer}, the real neighbourhoods behind Fig.~\ref{fig:gbpreal},
the meta-complex error of Fig.~\ref{fig:mergeerror}, the propagation cores of
Figs.~\ref{fig:corebonds} and~\ref{fig:coremeta}, and the branching sweep of
Fig.~\ref{fig:coretransition}. Those eleven are the only places
where the book quotes a number it does not recompute at build time, and each
says so where it is used. The real-network probe needs graph-tool and a network
connection the first time it runs, to pull its networks from Netzschleuder; it
caches their edge lists beside the other dumps the book uses, and after that
needs neither.

\section*{How the pieces fit}

Three points about the arrangement carry physical content rather than
bookkeeping, and they are the same three the construction was built around.

The intra-complex sum, Eq.~\eqref{eq:emitted}, takes the interior Hamiltonian as
an \emph{argument} rather than hard-coding it. That is what made
Sec.~\ref{sec:unanimity} a substitution rather than a rederivation, and it is
why a reader can add a model without touching anything else.

The hard-field and soft-field hitting-set solvers are kept apart rather than
unified behind one interface, because they answer different questions and agree
only at cardinality two. The routine that reports whether the hard-field cover
size can be believed returns true exactly when every cardinality is two ---
which, by Sec.~\ref{sec:corefree}, is exactly when the leaf-removal core has a
core-free branch. Two chapters' worth of argument is enforced in one predicate.

And the routine that returns a critical coupling \emph{refuses} the unanimity
interaction and directs the caller to the spinodal instead. For that model the
transition is usually discontinuous, and returning a spinodal under the name of
a critical coupling would be wrong by ten to fifty per cent, as
Sec.~\ref{sec:unanimity} shows. Enforcing this in the interface is more
reliable than a note in the documentation.

\section*{Equation to method}

Table~\ref{tab:map} is the map. Where a chapter has no manuscript behind it the
``evaluated by'' column names a function in the book's own \code{figs/} script,
and the check is in the same file.

{\scriptsize\setlength{\tabcolsep}{3pt}
\begin{longtable}{@{}L{0.08\textwidth} L{0.24\textwidth} L{0.36\textwidth} L{0.18\textwidth}@{}}
\caption{\textbf{Every computed equation, and what checks it.} Routines are in
\code{percolation/} for Part~II and \code{statmech/} for Parts~III and~IV
unless a path is given. Tests are the packages' own suites; \code{figs/} entries are the
book's scripts, which run their checks on every invocation.}
\label{tab:map}\\
\hline\hline
Eq. & what it is & evaluated by & checked in\\
\hline
\endfirsthead
\hline\hline
Eq. & what it is & evaluated by & checked in\\
\hline
\endhead
\hline\hline
\endfoot
\multicolumn{4}{@{}l}{\rule{0pt}{2.9ex}\textit{Chapter~\ref{ch:percolation}, percolation}}\\
\eqref{eq:Qminus} & the map, arriving from below & \code{giant.\allowbreak Chygraph.\allowbreak apply} & \code{test\_\allowbreak giant}\\
\eqref{eq:Qplus} & the map, arriving from above & \code{giant.\allowbreak Chygraph.\allowbreak apply} & \code{test\_\allowbreak giant}\\
\eqref{eq:Sl} & the order parameter & \code{giant.\allowbreak Chygraph.\allowbreak node\_\allowbreak fraction} & \code{test\_\allowbreak giant}\\
\eqref{eq:thetaH} & hypergraphs & \code{percolation.\allowbreak HypergraphPercolation} & \code{test\_\allowbreak percolation}\\
\eqref{eq:thetaG} & the graph threshold & \code{giant.\allowbreak Chygraph.\allowbreak theta} & \code{test\_\allowbreak percolation}\\
\eqref{eq:thetaMP} & multiplex networks & \code{percolation.\allowbreak MultiplexHypergraph} & \code{test\_\allowbreak percolation}\\
\eqref{eq:thetaNM} & links and triangles & \code{percolation.\allowbreak GraphWithTriangles} & \code{test\_\allowbreak percolation}\\
\eqref{eq:qc} & triangles at fixed degree & \code{figs/\allowbreak triangles\_\allowbreak vs\_\allowbreak regular.\allowbreak py} & same file\\
\multicolumn{4}{@{}l}{\rule{0pt}{2.9ex}\textit{Chapter~\ref{ch:giant}, beyond the threshold}}\\
\eqref{eq:AisJ} & the tensor is a Jacobian & \code{giant.\allowbreak Chygraph.\allowbreak jacobian}, \code{.\allowbreak A} & \code{test\_\allowbreak amplitude}\\
\eqref{eq:tripgf} & the bond-percolated triangle & \code{figs/\allowbreak triangles\_\allowbreak vs\_\allowbreak regular.\allowbreak py} & same file\\
\eqref{eq:Bgraph} & the critical amplitude & \code{amplitude.\allowbreak CriticalAmplitude} & \code{test\_\allowbreak amplitude}\\
\eqref{eq:Bhyper} & the same, hypergraphs & \code{amplitude.\allowbreak CriticalAmplitude} & \code{test\_\allowbreak amplitude}\\
\eqref{eq:jointmap} & correlated layers & \code{joint.\allowbreak JointGiantComponent} & \code{test\_\allowbreak joint}\\
\eqref{eq:biased} & inclusion-biased moments & \code{joint.\allowbreak JointGiantComponent} & \code{test\_\allowbreak joint}\\
\eqref{eq:qc3} & identical marginals, three thresholds & \code{joint.\allowbreak JointGiantComponent} & \code{figs/\allowbreak beyond\_\allowbreak threshold.\allowbreak py}\\
\eqref{eq:andor} & AND- against OR-logic & \code{applications.\allowbreak and\_\allowbreak or\_\allowbreak hypergraph} & \code{test\_\allowbreak applications}\\
\eqref{eq:stc} & triadic closure & \code{applications.\allowbreak stc\_\allowbreak percolation} & \code{test\_\allowbreak applications}\\
\eqref{eq:stcthresh} & its threshold, against Cirigliano & \code{applications.\allowbreak stc\_\allowbreak percolation} & \code{test\_\allowbreak applications}\\
\multicolumn{4}{@{}l}{\rule{0pt}{2.9ex}\textit{Chapter~\ref{ch:epidemics}, epidemics}}\\
\eqref{eq:Kn} & the household outbreak & \code{applications.\allowbreak clique\_\allowbreak excess\_\allowbreak pgf} & \code{figs/\allowbreak epidemics.\allowbreak py}\\
\eqref{eq:Rstar} & the household $R^{*}$ & \code{applications.\allowbreak household\_\allowbreak epidemic} & \code{figs/\allowbreak epidemics.\allowbreak py}\\
\eqref{eq:threshtau} & the threshold interior sum & \code{figs/\allowbreak epidemics.\allowbreak py:\allowbreak threshold\_\allowbreak recursion} & same file\\
\eqref{eq:threshmap} & the threshold-rule recursion & \code{figs/\allowbreak epidemics.\allowbreak py:\allowbreak check\_\allowbreak threshold\_\allowbreak recursion} & same file\\
\eqref{eq:contagion} & its Poisson two-layer case & \code{figs/\allowbreak epidemics.\allowbreak py:\allowbreak branch} & same file\\
\eqref{eq:tricritical} & its tricritical point & \code{figs/\allowbreak epidemics.\allowbreak py} & same file\\
\multicolumn{4}{@{}l}{\rule{0pt}{2.9ex}\textit{Chapter~\ref{ch:potts}, Potts}}\\
\eqref{eq:qonelimit} & the $q\to1$ limit & \code{figs/\allowbreak potts.\allowbreak py:\allowbreak percolation\_\allowbreak limit} & same file\\
\eqref{eq:tau2} & the link transmission & \code{figs/\allowbreak potts.\allowbreak py:\allowbreak transmission} & same file\\
\eqref{eq:tau3} & the triangle transmission & \code{figs/\allowbreak potts.\allowbreak py:\allowbreak transmission} & same file\\
\eqref{eq:branchcond} & one layer, both thresholds & \code{figs/\allowbreak potts.\allowbreak py:\allowbreak check\allowbreak \_thresholds} & same file\\
\eqref{eq:fold} & the ordered spinodal & \code{figs/\allowbreak potts.\allowbreak py:\allowbreak check\_\allowbreak transition\_\allowbreak point} & same file\\
\multicolumn{4}{@{}l}{\rule{0pt}{2.9ex}\textit{Chapter~\ref{ch:statmech}, the general theory}}\\
\eqref{eq:conv} & a generating function on a measure & \code{population.\allowbreak CavityPopulation.\allowbreak sweep} & \code{test\_\allowbreak population}\\
\eqref{eq:ensemblemap} & the ensemble map & \code{population.\allowbreak CavityPopulation.\allowbreak sweep} & \code{test\_\allowbreak population}\\
\eqref{eq:up} & the up step & \code{population.\allowbreak CavityPopulation.\allowbreak sweep} & \code{test\_\allowbreak population}\\
\eqref{eq:emitted} & the interior sum & \code{cavity.\allowbreak emitted\_\allowbreak field} & \code{test\_\allowbreak ising}\\
\eqref{eq:uprime} & $u'$ & \code{cavity.\allowbreak cavity\_\allowbreak derivative} & \code{test\_\allowbreak ising}\\
\eqref{eq:reweight} & the reweighting & \code{api.\allowbreak Chygraph.\allowbreak stability\_\allowbreak tensor} & \code{test\_\allowbreak stability}\\
\eqref{eq:branch} & the branching matrix & \code{api.\allowbreak Chygraph.\allowbreak branching\_\allowbreak matrix} & \code{test\_\allowbreak api}\\
\eqref{eq:sbaru} & the size-biased average & \code{api.\allowbreak Chygraph.\allowbreak transmission} & \code{test\_\allowbreak api}\\
\eqref{eq:det} & $\det(I-B)=0$ & \code{api.\allowbreak Chygraph.\allowbreak critical\_\allowbreak coupling} & \code{test\_\allowbreak api}\\
\eqref{eq:LT} & links and triangles & \code{api.\allowbreak Chygraph.\allowbreak critical\_\allowbreak coupling} & \code{test\_\allowbreak derivations}\\
\eqref{eq:exchange} & the exchange of stability & \code{fixedpoint.\allowbreak FixedPointStability} & \code{test\_\allowbreak fixedpoint}\\
\eqref{eq:bethe} & the Bethe free energy & \code{freeenergy.\allowbreak BetheFreeEnergy} & \code{test\_\allowbreak freeenergy}\\
\eqref{eq:para} & its paramagnetic closed form & \code{freeenergy.\allowbreak paramagnetic} & \code{test\_\allowbreak freeenergy}\\
\eqref{eq:sigmagen} & the complexity at $m=0$ & \code{figs/\allowbreak colouring.\allowbreak py:\allowbreak \_sv\_\allowbreak sigma} & same file\\
\multicolumn{4}{@{}l}{\rule{0pt}{2.9ex}\textit{Chapter~\ref{ch:ising}, the Ising model}}\\
\eqref{eq:Hising} & the Hamiltonian & \code{cavity.\allowbreak ising\_\allowbreak clique} & \code{test\_\allowbreak ising}\\
\eqref{eq:uclique} & $u'$ for $c=2,3,4$ & \code{ising.\allowbreak clique\_\allowbreak derivative} & \code{test\_\allowbreak derivations}\\
\eqref{eq:tcgraph} & the graph $T_{c}$ & \code{api.\allowbreak Chygraph.\allowbreak critical\_\allowbreak coupling} & \code{test\_\allowbreak ising}\\
\eqref{eq:mc} & Monte Carlo $T_{c}$ & \code{statmech/\allowbreak probe/\allowbreak ising\_\allowbreak mc.\allowbreak py} & --- (simulation)\\
\eqref{eq:simH} & the unanimity rule & \code{simplicial.\allowbreak emitted} & \code{test\_\allowbreak simplicial}\\
\eqref{eq:BW} & its infinite-connectivity limit & \code{simplicial.\allowbreak SimplicialChygraph.\allowbreak spinodal} & \code{test\_\allowbreak derivations}\\
\eqref{eq:Cq} & the residual $C_{q}$ & \code{simplicial.\allowbreak SimplicialChygraph.\allowbreak spinodal} & \code{figs/\allowbreak ising.\allowbreak py}\\
\multicolumn{4}{@{}l}{\rule{0pt}{2.9ex}\textit{Chapter~\ref{ch:hittingset}, hitting set}}\\
\eqref{eq:Hhs} & the constraint & \code{hittingset.\allowbreak HittingSet} & \code{test\_\allowbreak hittingset}\\
\eqref{eq:zhs} & warning propagation & \code{hittingset.\allowbreak HittingSet.\allowbreak F} & \code{test\_\allowbreak hittingset}\\
\eqref{eq:hs} & the hard-field map & \code{hittingset.\allowbreak HittingSet.\allowbreak solve} & \code{test\_\allowbreak hittingset}\\
\eqref{eq:kc} & $\ave{k}(c-1)=e$ & \code{hittingset.\allowbreak rsb\_\allowbreak point} & \code{test\_\allowbreak derivations}\\
\eqref{eq:soft} & soft-field BP & \code{softfield.\allowbreak HittingSetBP.\allowbreak sweep} & \code{test\_\allowbreak softfield}\\
\eqref{eq:hRS} & the regular solution & \code{softfield.\allowbreak regular\_\allowbreak field} & \code{test\_\allowbreak derivations}\\
\eqref{eq:sgen} & the entropy & \code{softfield.\allowbreak HittingSetBP.\allowbreak entropy} & \code{test\_\allowbreak softfield}\\
\eqref{eq:entropy} & its regular closed form & \code{softfield.\allowbreak regular\_\allowbreak entropy} & \code{test\_\allowbreak derivations}\\
\multicolumn{4}{@{}l}{\rule{0pt}{2.9ex}\textit{Chapter~\ref{ch:cover}, cover and core}}\\
\eqref{eq:Hvc} & the tighter constraint & \code{cover.\allowbreak CliqueCover} & \code{test\_\allowbreak cover}\\
\eqref{eq:vcz} & the $\max$ rule inside a complex & \code{cover.\allowbreak CliqueCover.\allowbreak tau} & \code{test\_\allowbreak cover}\\
\eqref{eq:vc} & the cover map & \code{cover.\allowbreak CliqueCover.\allowbreak solve} & \code{test\_\allowbreak cover}\\
\eqref{eq:core} & three-state leaf removal & \code{core.\allowbreak CorePercolation.\allowbreak solve} & \code{test\_\allowbreak core}\\
\eqref{eq:factored} & no core-free branch above $c=2$ & \code{core.\allowbreak CorePercolation.\allowbreak has\_\allowbreak core\_\allowbreak free\_\allowbreak branch} & \code{figs/\allowbreak cover.\allowbreak py}\\
\ref{tab:realcore} & leaf removal on real networks & \code{statmech/\allowbreak probe/\allowbreak real\_\allowbreak core.\allowbreak py} & \code{figs/\allowbreak cover.\allowbreak py:\allowbreak table\_\allowbreak real\_\allowbreak core}\\
\multicolumn{4}{@{}l}{\rule{0pt}{2.9ex}\textit{Chapter~\ref{ch:colouring}, colouring}}\\
\eqref{eq:taucol} & proper colouring, $-1/(q-1)$ & \code{figs/\allowbreak colouring.\allowbreak py:\allowbreak tau\_\allowbreak closed} & same file\\
\eqref{eq:tauhyp} & hypergraph colouring & \code{figs/\allowbreak colouring.\allowbreak py:\allowbreak tau\_\allowbreak closed} & same file\\
\eqref{eq:colthresh} & the stability threshold & \code{figs/\allowbreak colouring.\allowbreak py:\allowbreak stability\_\allowbreak threshold} & same file\\
\eqref{eq:coltwo} & proper against hypergraph & \code{figs/\allowbreak colouring.\allowbreak py:\allowbreak check\_\allowbreak cardinality} & same file\\
\eqref{eq:tridisagree} & triangles against the graph reading & \code{figs/\allowbreak colouring.\allowbreak py:\allowbreak triangle\_\allowbreak family} & same file\\
\eqref{eq:trifamily} & the family at fixed degree & \code{figs/\allowbreak colouring.\allowbreak py:\allowbreak check\_\allowbreak triangles\_\allowbreak at\_\allowbreak fixed\_\allowbreak degree} & same file\\
\eqref{eq:colsp} & the survey update & \code{figs/\allowbreak colouring.\allowbreak py:\allowbreak sp\_\allowbreak update} & same file\\
\eqref{eq:colsigma} & its complexity, against $c_{q}$ & \code{figs/\allowbreak colouring.\allowbreak py:\allowbreak sp\_\allowbreak threshold} & same file\\
\eqref{eq:hypsp} & cardinality, above two & \code{figs/\allowbreak colouring.\allowbreak py:\allowbreak hyp\_\allowbreak converge} & same file\\
\ref{fig:threshold} & the window, and where it closes & \code{figs/\allowbreak colouring.\allowbreak py:\allowbreak check\_\allowbreak window\_\allowbreak closes} & same file\\
\ref{tab:cq} & where colourings run out & \code{figs/\allowbreak colouring.\allowbreak py:\allowbreak check\_\allowbreak clustered\_\allowbreak cq} & same file\\
\multicolumn{4}{@{}l}{\rule{0pt}{2.9ex}\textit{Chapter~\ref{ch:satisfiability}, satisfiability}}\\
\eqref{eq:satZ} & the clause interior & \code{figs/\allowbreak satisfiability.\allowbreak py:\allowbreak clause\_\allowbreak Z} & same file\\
\eqref{eq:wp} & the warning map & \code{figs/\allowbreak satisfiability.\allowbreak py:\allowbreak wp\_\allowbreak map} & same file\\
\eqref{eq:satslope} & $\alpha=1$ at $k=2$, zero above & \code{figs/\allowbreak satisfiability.\allowbreak py} & same file\\
\eqref{eq:nested} & a clause whose member is a clause & \code{figs/\allowbreak satisfiability.\allowbreak py:\allowbreak nested\allowbreak \_G} & same file\\
\eqref{eq:flatten} & what CNF flattening costs & \code{figs/\allowbreak satisfiability.\allowbreak py:\allowbreak check\allowbreak \_nested\allowbreak \_flatten} & same file\\
\eqref{eq:sp} & the survey update & \code{figs/\allowbreak satisfiability.\allowbreak py:\allowbreak sp\_\allowbreak cavity\_\allowbreak piu} & same file\\
\eqref{eq:sigma} & its complexity, against $\alpha_{s}$ & \code{figs/\allowbreak satisfiability.\allowbreak py:\allowbreak sp\_\allowbreak threshold} & same file\\
\eqref{eq:nestedwp} & its warning map & \code{figs/\allowbreak satisfiability.\allowbreak py:\allowbreak nested\allowbreak \_G} & same file\\
\multicolumn{4}{@{}l}{\rule{0pt}{2.9ex}\textit{Chapter~\ref{ch:overlap}, overlapping complexes}}\\
--- & the cavity recursion on a chygraph & \code{statmech/\allowbreak probe/\allowbreak cavity\_\allowbreak clique.\allowbreak Chygraph\allowbreak BP} & \code{validate} in that file\\
\multicolumn{4}{@{}l}{\rule{0pt}{2.9ex}\textit{Chapter~\ref{ch:gbp}, generalised belief propagation}}\\
\eqref{eq:mobius} & M\"obius counting numbers & \code{region.\allowbreak RegionGraph} & \code{test\_\allowbreak region}\\
\eqref{eq:p2c} & the parent-to-child update & \code{gbp.\allowbreak GBP} & \code{test\_\allowbreak gbp}\\
\multicolumn{4}{@{}l}{\rule{0pt}{2.9ex}\textit{Chapter~\ref{ch:metacomplex}, meta-complexes}}\\
\eqref{eq:branching} & the incidence branching & \code{statmech/\allowbreak probe/\allowbreak karrer\_\allowbreak core\_\allowbreak sweep.\allowbreak py} & \code{check\_\allowbreak core\_\allowbreak transition}\\
\eqref{eq:corea} & the pruning fixed point & \code{statmech/\allowbreak probe/\allowbreak karrer\_\allowbreak core\_\allowbreak sweep.\allowbreak core\_\allowbreak closed\_\allowbreak form} & \code{check\_\allowbreak closed\_\allowbreak form}\\
\eqref{eq:corefrac} & the core fraction & \code{statmech/\allowbreak probe/\allowbreak karrer\_\allowbreak core\_\allowbreak sweep.\allowbreak core\_\allowbreak closed\_\allowbreak form} & \code{check\_\allowbreak closed\_\allowbreak form}\\
\end{longtable}
}

Two entries need a word. Equation~\eqref{eq:emitted} is a placeholder for
whatever the interior Hamiltonian is: the routine takes $-\beta H$ as a
callable, and the models of Chapters~\ref{ch:ising}--\ref{ch:satisfiability}
differ only in what is passed to it. And Eq.~\eqref{eq:mc} is a direct
simulation, deliberately not covered by the test suite --- a test that
reproduced it from the cavity equations would defeat its purpose. It is the one
number in the book obtained without them.

The remaining routines are solvers shared across chapters rather than single
equations. \code{antimonotone.\allowbreak fixed\_\allowbreak point} brackets the order-reversing maps of
Chapters~\ref{ch:hittingset} and~\ref{ch:cover} by iterating $F\circ F$ from $0$
and from $1$, which is the construction of Sec.~\ref{sec:hard};
\code{region.\allowbreak overlap\_\allowbreak profile} measures how far an explicit complex family is
from treelike; and \code{stability.\allowbreak StabilityMatrix} is the symbolic $2L^{2}$
tensor of Sec.~\ref{sec:threshold} with the reweighting of
Eq.~\eqref{eq:reweight} applied.

\section*{Reproducing a figure}

Every figure in the book that is not drawn in TikZ is generated by a script in
\code{figs/\allowbreak }, usually one per chapter and named for it. Running a
script rewrites its PDFs into the same
directory and prints the checks it performs, so
\begin{quote}
\code{python figs/\allowbreak ising.\allowbreak py}
\end{quote}
regenerates Chapter~\ref{ch:ising}'s three figures and reports, among other
things, that $u'$ agrees with its closed form, that $T_{c}$ falls monotonically
with clustering at every degree, and where the unanimity transition stops being
continuous. The build then picks up the new PDFs with no further intervention.

Several of the scripts are slow, and for the same reason: they bisect a
threshold or a coexistence window, over more than one seed.
\code{figs/\allowbreak colouring.\allowbreak py} is the slowest in the book at
about fourteen minutes and \code{figs/\allowbreak satisfiability.\allowbreak py}
takes about five, both because of their survey-propagation sections;
\code{figs/\allowbreak merge.\allowbreak py},
\code{figs/\allowbreak ising.\allowbreak py} and
\code{figs/\allowbreak potts.\allowbreak py} take a few minutes each. The rest
run in seconds.
